%% arara directives
% arara: xelatex
% arara: bibtex
% arara: xelatex
% arara: xelatex

%\documentclass{article} % One-column default
\documentclass[twocolumn, switch]{article} % Method A for two-column formatting

\usepackage{preprint}

%% Math packages
\usepackage{amsmath, amsthm, amssymb, amsfonts, algorithm2e, algpseudocode, graphicx, caption, subcaption}
\RestyleAlgo{ruled}
\graphicspath{ {./Pictures/} }
%% Bibliography options
\usepackage[numbers,square]{natbib}
\bibliographystyle{unsrtnat}
%\usepackage{natbib}
%\bibliographystyle{Geology}

%% General packages
\usepackage[utf8]{inputenc}	% allow utf-8 input
\usepackage[T1]{fontenc}	% use 8-bit T1 fonts
\usepackage{xcolor}		% colors for hyperlinks
\usepackage[colorlinks = true,
            linkcolor = purple,
            urlcolor  = blue,
            citecolor = cyan,
            anchorcolor = black]{hyperref}	% Color links to references, figures, etc.
\usepackage{booktabs} 		% professional-quality tables
\usepackage{nicefrac}		% compact symbols for 1/2, etc.
\usepackage{microtype}		% microtypography
\usepackage{lineno}		% Line numbers
\usepackage{float}			% Allows for figures within multicol
%\usepackage{multicol}		% Multiple columns (Method B)

\usepackage{lipsum}		%  Filler text

 %% Special figure caption options
\usepackage{newfloat}
\DeclareFloatingEnvironment[name={Supplementary Figure}]{suppfigure}
\usepackage{sidecap}
\sidecaptionvpos{figure}{c}

% Section title spacing  options
\usepackage{titlesec}
\titlespacing\section{0pt}{12pt plus 3pt minus 3pt}{1pt plus 1pt minus 1pt}
\titlespacing\subsection{0pt}{10pt plus 3pt minus 3pt}{1pt plus 1pt minus 1pt}
\titlespacing\subsubsection{0pt}{8pt plus 3pt minus 3pt}{1pt plus 1pt minus 1pt}

% ORCiD insertion
\usepackage{tikz,xcolor,hyperref}

\definecolor{lime}{HTML}{A6CE39}
\DeclareRobustCommand{\orcidicon}{
	\begin{tikzpicture}
	\draw[lime, fill=lime] (0,0) 
	circle [radius=0.16] 
	node[white] {{\fontfamily{qag}\selectfont \tiny ID}};
	\draw[white, fill=white] (-0.0625,0.095) 
	circle [radius=0.007];
	\end{tikzpicture}
	\hspace{-2mm}
}
\foreach \x in {A, ..., Z}{\expandafter\xdef\csname orcid\x\endcsname{\noexpand\href{https://orcid.org/\csname orcidauthor\x\endcsname}
			{\noexpand\orcidicon}}
}
% Define the ORCID iD command for each author separately. Here done for two authors.
\newcommand{\orcidauthorA}{0000-0000-0000-0001}
\newcommand{\orcidauthorB}{0000-0000-0000-0002}
\newcommand{\orcidauthorC}{0000-0000-0000-0003}
\newcommand{\orcidauthorD}{0000-0000-0000-0004}

%%%%%%%%%%%%%%%%   Title   %%%%%%%%%%%%%%%%
\title{Parallel A* project}

% Add watermark with submission status
\usepackage{xwatermark}
% Left watermark
\newwatermark[firstpage,color=gray!60,angle=90,scale=0.32, xpos=-4.05in,ypos=0]{\href{https://doi.org/}{\color{gray}{Publication doi}}}
% Right watermark
\newwatermark[firstpage,color=gray!60,angle=90,scale=0.32, xpos=3.9in,ypos=0]{\href{https://doi.org/}{\color{gray}{Preprint doi}}}
% Bottom watermark
\newwatermark[firstpage,color=gray!90,angle=0,scale=0.28, xpos=0in,ypos=-5in]{*correspondence: \texttt{email@institution.edu}}

%%%%%%%%%%%%%%%  Author list  %%%%%%%%%%%%%%%
\usepackage{authblk}
\renewcommand*{\Authfont}{\bfseries}
\author{Lorenzo Ippolito, Mattia Rosso, Fabio Mirto}


%%%%%%%%%%%%%%    Front matter    %%%%%%%%%%%%%%
\begin{document}

\twocolumn[ % Method A for two-column formatting
  \begin{@twocolumnfalse} % Method A for two-column formatting
  
\maketitle

\begin{abstract}
  ...Here the abstract...
\end{abstract}
%\keywords{First keyword \and Second keyword \and More} % (optional)
\vspace{0.35cm}

  \end{@twocolumnfalse} % Method A for two-column formatting
] % Method A for two-column formatting

%\begin{multicols}{2} % Method B for two-column formatting (doesn't play well with line numbers), comment out if using method A


%%%%%%%%%%%%%%%  Main text   %%%%%%%%%%%%%%%
% \linenumbers


\section{Introduction: about the A* algorithm}
(Explanation of the general A* problem)

\section{A* project application}
(What we will apply the algorithm to)

\subsection{Heuristic function}
(Explanation + formula)

\section{Test data and benchmark}
(Explanation, include file reading problem)
\paragraph{Random Graph}
(Explanation about how the random graph was generated)
\paragraph{Benchmark}
(General details about the benchmark used and about the paths chosed)
\begin{figure}[ht!]
  \centering
  \includegraphics[width=0.5\linewidth]{google_maps.png}
  \caption{A subfigure}
  \label{fig:sub2}
\end{figure}

\section{Sequential A* Algorithm}
(Explanation)

\subsection{Psedocode}
\SetKwComment{Comment}{/* }{ */}
\begin{algorithm} 
\caption{An algorithm with caption}\label{alg:two}
\KwData{$n \geq 0$}
\KwResult{$y = x^n$}
$y \gets 1$\;
$X \gets x$\;
$N \gets n$\;
\While{$N \neq 0$}{
  \eIf{$N$ is even}{
    $X \gets X \times X$\;
    $N \gets \frac{N}{2}$ \Comment*[r]{This is a comment}
  }{\If{$N$ is odd}{
      $y \gets y \times X$\;
      $N \gets N - 1$\;
    }
  }
}
\end{algorithm}

\subsection{C Implementation}
(Explain the details of the input file and of the main data structures used)

\subsection{Results}
(Results of sequential reading + A* sequential on random, BAY, FLA)
(Table with numbers)

\section{A* and Dijkstra: a comparison}
(Explanation)

\subsection{Results}
(Picture A* vs Dijkstra on BAY and FLA)

\section{Parallel reading of the input file}
(Motivation of parallel reading)
\subsection{Parallel Read: approach 1}
(Explanation)
\subsection{Parallel Read: approach 2}
(Explanation)
\subsection{Parallel Read: approach 3}
(Explanation)
\subsection{Results}
(Plots of parallel reading vs sequential reading)

\section{Parallel A*: two examinated approaches}
(Explanation)
\subsection{First Attempt In Parallelizing A*}
(Explanation + Pseudocode?)
\subsection{HDA*}
(Introduction + map with colored points)
\subsubsection{Message Passing Model}
(Explanation + Pseudocode?)
\subsubsection{Shared Address Space Model}
\paragraph{Barrier(SAS-B)}
(Explanation + Pseudocode?)
\paragraph{Barrier(SAS-SF)}
(Explanation + Pseudocode?)
\subsection{Results}
(Plots with the comparison of all the models with sequential A* on
random, BAY, FLA)

\section{Complete Results}
(Tables with numbers)

\section{Final Considerations}
(Comments)

\section{DIMACS Benchmark}
(More detailed explanation of the input format of the benchmarks)

\section{Future Works}
(Possible improvements)




%%%%%%%%%%%% Supplementary Methods %%%%%%%%%%%%
%\footnotesize
%\section*{Methods}

%%%%%%%%%%%%% Acknowledgements %%%%%%%%%%%%%
%\footnotesize
%\section*{Acknowledgements}

%%%%%%%%%%%%%%   Bibliography   %%%%%%%%%%%%%%
\normalsize
\bibliography{references}

%%%%%%%%%%%%  Supplementary Figures  %%%%%%%%%%%%
%\clearpage

%%%%%%%%%%%%%%%%   End   %%%%%%%%%%%%%%%%
%\end{multicols}  % Method B for two-column formatting (doesn't play well with line numbers), comment out if using method A
\end{document}